% Generated draft: requires mathematical and source review.
\documentclass[11pt]{article}
\usepackage[margin=25mm]{geometry}
\usepackage{fontspec}
\setmainfont{Latin Modern Roman}
\usepackage{amsmath,amssymb,mathtools,unicode-math}
\setmathfont{Latin Modern Math}
\usepackage{enumitem,needspace,xurl,hyperref}
\setlength{\parindent}{0pt}
\setlength{\parskip}{5pt}
\setlength{\emergencystretch}{3em}
\newcommand{\N}{\mathbb{N}}
\newcommand{\Z}{\mathbb{Z}}
\newcommand{\Q}{\mathbb{Q}}
\newcommand{\R}{\mathbb{R}}
\newcommand{\C}{\mathbb{C}}
\newcommand{\F}{\mathbb{F}}
\newcommand{\A}{\mathbb{A}}
\newcommand{\PP}{\mathbb{P}}
\newcommand{\E}{\mathbb{E}}
\newcommand{\eps}{\varepsilon}
\newcommand{\dd}{\,\mathrm{d}}
\begin{document}
\section*{UnsolvedMath ID: 30005617}
\textbf{Source problem code:} catalogue entry
\subsection{Definitions and mathematical statement}
The supplied record combines two separate problems. Its waveguide formulation is reconstructed, while the quoted original passage begins with remarks referring to an omitted preceding question.

For the waveguide problem, let \(H\) denote the self-adjoint operator associated with an admissible curved waveguide generated by a curve \(\Gamma\), with confinement by a potential \(V\) in the soft case or by a Dirichlet boundary condition in the hard-wall case. The supplied record does not specify the differential expression, ambient dimension, operator domain, admissible geometry, or hypotheses on \(V\). Consequently, the admissible family of operators remains unresolved. Write \(\sigma_{\mathrm{disc}}(H)\) for the isolated eigenvalues of finite multiplicity. When \(H\) has a ground-state eigenvalue \(E_0=\inf\sigma(H)\), the question in the reconstructed statement is whether every admissible \(H\) satisfies
\[
\forall\lambda\in\sigma_{\mathrm{disc}}(H),\qquad
\lambda>E_0\ \Longrightarrow\ \dim\ker(H-\lambda I)=1.
\]
This expresses simplicity of each discrete eigenvalue strictly above the ground state; it does not assert simplicity of the ground state. A precise universal conjecture requires the missing definition of the admissible waveguides. The supplied excerpt states that suitable choices of \(\Gamma\) and \(V\) can produce any finite number of eigenvalues, that the question remains open for hard-wall waveguides, and that the quasi-one-dimensional restriction is essential because replacing the generating curve by a graph can yield nonsimple spectrum. These are supplied source qualifications, not independently verified assertions.

For the surface problem, fix an area \(A>0\) and an integer \(g\geq1\). Let \(\mathcal S_{A,g}\) denote the admissible connected compact smooth embedded surfaces \(\Sigma\subset\R^3\) with
\[
\operatorname{Area}(\Sigma)=A,\qquad \operatorname{genus}(\Sigma)=g.
\]
The source does not explicitly specify whether surfaces must have empty boundary. That convention, and any conditions on an allowed boundary, must be fixed before \(\mathcal S_{A,g}\) is fully defined. At a point \(p\in\Sigma\), let \(k_1(p),k_2(p)\) be the principal curvatures computed using a local unit normal, and define
\[
H_\Sigma:=\|k_1-k_2\|_{L^\infty(\Sigma)}
=\sup_{p\in\Sigma}|k_1(p)-k_2(p)|.
\]
The absolute difference is independent of the ordering of the principal curvatures and of reversal of the unit normal. Determine
\[
m(A,g):=\inf_{\Sigma\in\mathcal S_{A,g}}H_\Sigma,
\]
obtain lower bounds for this infimum, decide whether it is attained, and characterize all attaining surfaces if they exist. The excerpt singles out the torus case \(g=1\) as particularly interesting. It also says that \(H_\Sigma\) occurs as the ``size'' of the essential spectrum of Dirac operators with special transmission conditions along \(\Sigma\), referring to an unidentified reference [3]; neither those operators nor the meaning of ``size'' is supplied.
\subsection{Short English statement}
The reconstructed waveguide question asks whether every discrete eigenvalue above an existing ground state has multiplicity one, for both potential confinement and Dirichlet confinement. The admissible operators are not specified. Separately, for each fixed area and genus \(g\geq1\), minimize the maximum absolute difference of the principal curvatures and describe minimizers, if any exist. The surface boundary convention remains unresolved.
\subsection{Sources}
\begin{itemize}
\item Local UnsolvedMath registry, \texttt{data/unsolvedmath/problems.json}, permanent ID \(30005617\), source problem code \texttt{catalogue entry}, title ``Curved-Waveguide Simplicity and Surface-Curvature Optimization.'' This draft uses the supplied complete registry record; no independent registry or external-source inspection is claimed.
\item Supplied report citation: \emph{Geometric Spectral Theory} (2024), \emph{Oberwolfach Reports}, pp.\ 2017--2106, DOI \texttt{10.4171/owr/2023/36}. Supplied source URL and literature-source link, labelled ``original OWR report'': \url{https://doi.org/10.4171/owr/2023/36}. The background describes these reports as records of mathematical workshops and their open-problem discussions.
\item The registry records proposed year 2023 and supplies no proposer attribution. The citation labels the report 2024, while the DOI and year tag contain 2023; their relationship is not explained in the supplied record.
\item Literature assessment dated exactly \texttt{2026-08-22}: the supplied curation labels the problem open and states that the finite-eigenvalue question remains open even for hard-wall waveguides and that the cited literature provides no general resolution. This is supplied dated curation, not an independently verified literature assessment. Its wording does not unambiguously identify the reconstructed simplicity question.
\item The supplied statement status is \texttt{corrected\_verified}; its verification note says the statement was reconstructed and checked against the official Oberwolfach report. This is a supplied provenance claim, not verification performed for this draft.
\item The quoted passage cites [3] for the Dirac-operator connection, but no bibliographic entry for [3] is supplied. The report is the supplied authoritative citation; no separately identifiable reference for that connection is available in the record.
\end{itemize}
\end{document}
